% !TEX root = hw10-answers.tex
\chead{\textbf{Exercises week 10}}
\begin{document}
\flushleft\textbf{Topic: Structural Causal Models}

\begin{enumerate}
\item ($1 + 3\times 2$ points) Prove the following remarks on the equivalence of SCMs from the lecture:

\begin{enumerate}
  \item Prove that equivalence of SCMs ($M \equiv \tilde{M}$) is indeed an equivalence relation.  
  \item \label{ex:equiv} Equivalence is preserved by hard interventions, i.e.: if $M$ is equivalent to $\tilde{M}$ and $T \subseteq J \cup V \cup W$ is an intervention target, then $M_{\doit(T\dots)}$ is equivalent to $\tilde{M}_{\doit(T\dots)}$. Do this for the three variants of hard interventions.  
  \item Show that equivalent SCMs have the same solution functions (Definition \ref*{def:solution-function}). Deduce from this and from Exercise \ref{ex:equiv} that unique solvability and simplicity are invariants of equivalence, i.e.: If $M$ is equivalent to $\tilde{M}$ and $M$ is uniquely solvable or simple, then $\tilde{M}$ is also uniquely solvable or simple, respectively. 
  \item Equivalent SCMs have the same potential outcomes (Definition \ref*{def:scm_solutions_for_inputs}), solutions (Definition \ref*{def:scm_solutions}), and the same Markov kernels (Definition \ref*{def:scm_kernels}). 
\end{enumerate}


\begin{gradedsolution}
In all parts we assume that $M = \langle J, V, W, \mathcal{X}, P, f \rangle$ and $\tilde{M} = \langle J, V, W, \mathcal{X}, P, \tilde{f} \rangle$ are two SCMs which can only differ in their causal mechanisms.
\begin{enumerate}
	\item This follows directly from Definition \ref*{def:equivalent_scms} and the fact that logical equivalence ``$\Longleftrightarrow$'' is an equivalence relation. 
	\item We discuss only the case of hard interventions with specified target value, since the other two cases work the same way (and are even slightly easier to show).
	Thus, let $T \subseteq J \cup V \cup W$ be an intervention target and $\xi_{T} \in \mathcal{X}_T$ be the intervention target value. 
	The intervened SCMs are then given by 
	\begin{align*}
		M_{\doit(X_T = \xi_T)} = \left\langle J \setminus T, V \cup T, W \setminus T, \mathcal{X}, P_{W \setminus T}, (f_{V \setminus T}, \xi_T) \right\rangle, \\
		\tilde{M}_{\doit(X_T = \xi_T)} = \left\langle J \setminus T, V \cup T, W \setminus T, \mathcal{X}, P_{W \setminus T}, (\tilde{f}_{V \setminus T}, \xi_T) \right\rangle.
	\end{align*}
	Now, set $f' := (f_{V \setminus T}, \xi_T)$ and $\tilde{f}' := (\tilde{f}_{V \setminus T}, \xi_T)$. 

	Then, for $v \in V \setminus T$ and arbitrary $x \in \mathcal{X}$ we obtain:
	\begin{align*}
		x_v = f'_v(x) & \Longleftrightarrow x_v = f_v(x) \\
		& \Longleftrightarrow x_v = \tilde{f}_v(x) \\
		& \Longleftrightarrow x_v = \tilde{f}'_v(x).
	\end{align*}
	Thereby, in the first step we used the definition of $f'$ and that $v \in V \setminus T$. 
	In the second step, we used the equivalence $M \equiv \tilde{M}$.
	In the last step, we used the definition of $\tilde{f}'$ and again that $v \in V \setminus T$.

	Finally, for $v \in T$ and arbitrary $x \in X$, we obtain:
	\begin{align*}
		x_v = f'_v(x) & \Longleftrightarrow x_v = \xi_v \\
		& \Longleftrightarrow x_v = \tilde{f}'_v(x).
	\end{align*}
	Since $V \cup T = (V \setminus T) \cup T$, this proves the claim.
	\item Let $M \equiv \tilde{M}$, and let $g: \mathcal{X}_J \times \mathcal{X}_W \to \mathcal{X}_V$ be any measurable function.
	Then the following equivalences hold:
	\begin{align*}
		g \text{ is a solution function for } M & \Longleftrightarrow \forall x_J, \forall x_W: \ \ g(x_J, x_W) = f(x_J, g(x_J, x_W), x_W) \\
		&   \Longleftrightarrow \forall v, \forall x_J, \forall x_W: \ \ g_v(x_J, x_W) = f_v(x_J, g(x_J, x_W), x_W) \\
		&  \Longleftrightarrow   \forall v, \forall x_J, \forall x_W: \ \ g_v(x_J, x_W) = \tilde{f}_v(x_J, g(x_J, x_W), x_W) \\
		& \Longleftrightarrow   \forall x_J, \forall x_W: \ \ g(x_J, x_W) = \tilde{f}(x_J, g(x_J, x_W), x_W) \\
		& \Longleftrightarrow g \text{ is a solution function for } \tilde{M}.
	\end{align*}
	This shows that equivalent SCMs have the same solution functions, and consequently that unique solvability is an invariant of equivalence.

	Now, assume that $M$ is simple and that $M \equiv \tilde{M}$.
	We want to show that $\tilde{M}$ is also simple.
	For this, let $T \subseteq J \cup V \cup W$ be any subset.
	Then we need to show that $\tilde{M}_{\doit(T)}$ is uniquely solvable.
	But note that due to the simplicity of $M$, $M_{\doit(T)}$ is uniquely solvable. 
	Furthermore, as was shown in Exercise \ref{ex:equiv}, $M_{\doit(T)} \equiv \tilde{M}_{\doit(T)}$. We have just shown that unique solvability is an invariant of equivalence, and thus the unique solvability of $\tilde{M}_{\doit(T)}$ automatically follows.

	\item Let $M \equiv \tilde{M}$.
	Let $X_{V\cup W}^{\doit(x_J)}$ be any random variable with codomain $\mathcal{X}_V \times \mathcal{X}_W$, and with $x_J \in \mathcal{X}_J$.
	We have the following equivalences:
	\begin{align*}
		&	X_{V \cup W}^{\doit(x_J)} \text{ is a potential outcome of } M \text{ for input } x_J \in \mathcal{X}_J \\
		\Longleftrightarrow & X_W^{\doit(x_J)} \sim P \text{ and } X_V^{\doit(x_J)} = f(x_J, X_V^{\doit(x_J)}, X_W^{\doit(x_J)}) \text{ a.s. } \\
		\Longleftrightarrow & X_W^{\doit(x_J)} \sim P \text{ and } \forall v \in V: \ \ X_v^{\doit(x_J)} = f_v(x_J, X_V^{\doit(x_J)}, X_W^{\doit(x_J)}) \text{ a.s. } \\
		\Longleftrightarrow & X_W^{\doit(x_J)} \sim P \text{ and } \forall v \in V: \ \ X_v^{\doit(x_J)} = \tilde{f}_v(x_J, X_V^{\doit(x_J)}, X_W^{\doit(x_J)}) \text{ a.s. } \\
		\Longleftrightarrow & X_W^{\doit(x_J)} \sim P \text{ and } X_V^{\doit(x_J)} = \tilde{f}(x_J, X_V^{\doit(x_J)}, X_W^{\doit(x_J)}) \text{ a.s. } \\
		\Longleftrightarrow  &	X_{V \cup W}^{\doit(x_J)} \text{ is a solution for } \tilde{M} \text{ with input } x_J \in \mathcal{X}_J. 
	\end{align*}
	This shows that equivalent SCMs have the same potential outcomes. Since this holds for all $x_J \in \Xcal_J$ and potential outcomes are, by definition, just evaluations of a solution, we have that equivalent SCMs have the same solutions as well.

	From Definition \ref*{def:scm_kernels} it is immediate that SCMs with the same solutions have the same Markov kernels as well.
\end{enumerate}
\end{gradedsolution}

\item ($6 \times 1$ points) Given SCM: $M = \left\langle J=\emptyset, V=\{A, B, C, D\}, W=\{A', B', C', D'\}, \R^8, P, f \right\rangle$, where $P$ are four independent normal distributions with means and standard deviations $\mu_A, \sigma_A, \mu_B, ..., \sigma_D$, and
\[
\begin{pmatrix}
	f_A(x, \varepsilon) \\ f_B(x, \varepsilon) \\ f_C(x, \varepsilon) \\ f_D(x, \varepsilon)
\end{pmatrix}
= 
\begin{pmatrix}
	0 & 0 & 0 & 0 \\
	\alpha_{ab} & 0 & 0 & \alpha_{db} \\
	0 & \alpha_{bc} & 0 & \alpha_{dc} \\
	\alpha_{ad} & 0 & 0 & 0 \\
\end{pmatrix}
\begin{pmatrix}
	x_A \\ x_B \\ x_C \\ x_D
\end{pmatrix}
+
\begin{pmatrix}
	\varepsilon_{A'} \\ \varepsilon_{B'} \\ \varepsilon_{C'} \\ \varepsilon_{D'}
\end{pmatrix}
\]
for nonzero real numbers $\alpha_{(\cdot)}$. Note that we have a slightly different notation as in the lecture notes, as we refer to $x_V$ as $x$, and refer to $x_W$ as $\varepsilon$.
\begin{enumerate}
	\item Construct a solution function. 
	\item Give the graph $G(M)$. 
	\item For latent variables $L=\{D\}$, give the marginal SCM $M_{\setminus L}$ and its solution function. 
	\item Show $G(M_{\setminus L})$ is a subgraph of $G(M)^{\setminus L}$. 
	\item Show observational equivalence of $M$ and $M_{\setminus L}$ on $V \setminus L$ directly using the solution functions and Markov kernels. 
	\item Show observational equivalence between $M_{\doit(B)}$ and ${\left(M_{\setminus L}\right)}_{\doit(B)}$ on $\{A,C\}$. 
\end{enumerate}

\begin{gradedsolution}
\begin{enumerate}
	\item As the SCM is acyclic with order e.g. $(A,D,B,C)$, the solution can be read off from $f$. Alternatively, use  $(1_4-B_{VV})^{-1}\varepsilon$.
	One finds:
	\[\begin{cases}
		g_A(\varepsilon) = \varepsilon_{A'} \\
		g_B(\varepsilon) = \varepsilon_{A'} \left(\alpha _{ab}+\alpha _{ad} \alpha _{db}\right)+\varepsilon_{B'}+\varepsilon_{D'} \alpha _{db} \\
		g_C(\varepsilon) = \varepsilon_{A'} \left(\alpha _{ab} \alpha _{bc}+\alpha _{ad} \alpha _{bc} \alpha _{db}+\alpha _{ad} \alpha _{dc}\right)+\varepsilon_{B'} \alpha _{bc}+\varepsilon_{D'} \left(\alpha _{bc} \alpha _{db}+\alpha _{dc}\right)+\varepsilon_{C'} \\
		g_D(\varepsilon) = \varepsilon_{A'} \alpha _{ad}+\varepsilon_{D'}
	\end{cases}\]
	\item $G(M):$
	\begin{center}
		\begin{tikzpicture}[scale=1, transform shape]
			\tikzstyle{every node} = [draw,shape=circle];
			\node (A) at (0, 2) {$A$};
			\node (B) at (1, 0) {$B$};
			\node (C) at (3, 0) {$C$};
			\node (D) at (2, 2) {$D$};
			\node (epsA) at (0, 4) {${A'}$};
			\node (epsB) at (1, -2) {${B'}$};
			\node (epsC) at (3, -2) {${C'}$};
			\node (epsD) at (2, 4) {${D'}$};
			\draw[-{Latex[length=3mm, width=2mm]}] (A) to (B);
			\draw[-{Latex[length=3mm, width=2mm]}] (A) to (D);
			\draw[-{Latex[length=3mm, width=2mm]}] (D) to (C);
			\draw[-{Latex[length=3mm, width=2mm]}] (D) to (B);
			\draw[-{Latex[length=3mm, width=2mm]}] (B) to (C);
			\draw[-{Latex[length=3mm, width=2mm]}] (epsA) to (A);
			\draw[-{Latex[length=3mm, width=2mm]}] (epsB) to (B);
			\draw[-{Latex[length=3mm, width=2mm]}] (epsC) to (C);
			\draw[-{Latex[length=3mm, width=2mm]}] (epsD) to (D);
		\end{tikzpicture}
	\end{center}
		
	\item Intervening on $\{A,B,C\}$, we solve for $x_D$ and find $g'_D(x_{(A, B, C)}, \varepsilon)=\alpha_{ad} x_A + \varepsilon_{D'}$ as a solution function of $M_{\doit(A, B, C)}$. Filling this in at $f$ leads to:
	\begin{align*}
		\begin{pmatrix}
			\tilde{f}_A(x) \\ \tilde{f}_B(x) \\ \tilde{f}_C(x)
		\end{pmatrix}
		&= 
		\begin{pmatrix}
			0 & 0 & 0 & 0 \\
			\alpha_{ab} & 0 & 0 & \alpha_{db} \\
			0 & \alpha_{bc} & 0 & \alpha_{dc} \\
		\end{pmatrix}
		\begin{pmatrix}
			x_A \\ x_B \\ x_C \\ \alpha_{ad} x_A + \varepsilon_{D'}
		\end{pmatrix}
		+
		\begin{pmatrix}
			\varepsilon_{A'} \\ \varepsilon_{B'} \\ \varepsilon_{C'}
		\end{pmatrix} \\
		&= 
		\begin{pmatrix}
			0 & 0 & 0 \\
			\alpha_{ab}+\alpha_{ad}\alpha_{db} & 0 & 0 \\
			\alpha_{ad}\alpha_{dc} & \alpha_{bc} & 0 \\
		\end{pmatrix}
		\begin{pmatrix}
			x_A \\ x_B \\ x_C
		\end{pmatrix}
		+
		\begin{pmatrix}
			\varepsilon_{A'} \\ \varepsilon_{B'} +\alpha_{db}\varepsilon_{D'}\\ \varepsilon_{C'} + \alpha_{dc}\varepsilon_{D'}
		\end{pmatrix}
	\end{align*}
	as the causal mechanism of $M_{\setminus L}$. The solution function of $M_{\setminus L}$ is
	\[\begin{cases}
		\tilde{g}_A(\varepsilon) = \varepsilon_{A'} \\
		\tilde{g}_B(\varepsilon) = \varepsilon_{A'} \left(\alpha _{ab}+\alpha _{ad} \alpha _{db}\right)+\varepsilon_{B'}+\varepsilon_{D'} \alpha _{db} \\
		\tilde{g}_C(\varepsilon) = \varepsilon_{A'} \left(\alpha _{ab} \alpha _{bc}+\alpha _{ad} \alpha _{bc} \alpha _{db}+\alpha _{ad} \alpha _{dc}\right)+\varepsilon_{B'} \alpha _{bc}+\varepsilon_{D'} \left(\alpha _{bc} \alpha _{db}+\alpha _{dc}\right)+\varepsilon_{C'},
	\end{cases}\]
	which follows directly from Lemma \ref*{lem:marginalization_preserves_unique_solvability}.
	
	\item $G(M_{\setminus L}):$
	 \begin{center}
		\begin{tikzpicture}[scale=1, transform shape]
			\tikzstyle{every node} = [draw,shape=circle];
			\node (A) at (0, 2) {$A$};
			\node (B) at (1, 0) {$B$};
			\node (C) at (3, 0) {$C$};
			\node (epsA) at (0, 4) {${A'}$};
			\node (epsB) at (1, -2) {${B'}$};
			\node (epsC) at (3, -2) {${C'}$};
			\node (epsD) at (2, 4) {${D'}$};
			\draw[-{Latex[length=3mm, width=2mm]}] (A) to (B);
			\draw[-{Latex[length=3mm, width=2mm]}] (A) to (B);
			\draw[-{Latex[length=3mm, width=2mm]}] (A) to (C);
			\draw[-{Latex[length=3mm, width=2mm]}] (epsD) to (B);
			\draw[-{Latex[length=3mm, width=2mm]}] (epsD) to (C);
			\draw[-{Latex[length=3mm, width=2mm]}] (B) to (C);
			\draw[-{Latex[length=3mm, width=2mm]}] (epsA) to (A);
			\draw[-{Latex[length=3mm, width=2mm]}] (epsB) to (B);
			\draw[-{Latex[length=3mm, width=2mm]}] (epsC) to (C);
		\end{tikzpicture}
	\end{center}
	
	$G(M)^{\setminus L}:$
	\begin{center}
		\begin{tikzpicture}[scale=1, transform shape]
			\tikzstyle{every node} = [draw,shape=circle];
			\node (A) at (0, 2) {$A$};
			\node (B) at (1, 0) {$B$};
			\node (C) at (3, 0) {$C$};
			\node (epsA) at (0, 4) {${A'}$};
			\node (epsB) at (1, -2) {${B'}$};
			\node (epsC) at (3, -2) {${C'}$};
			\node (epsD) at (2, 4) {${D'}$};
			\draw[-{Latex[length=3mm, width=2mm]}] (A) to (B);
			\draw[-{Latex[length=3mm, width=2mm]}] (A) to (B);
			\draw[-{Latex[length=3mm, width=2mm]}] (A) to (C);
			\draw[-{Latex[length=3mm, width=2mm]}] (epsD) to (B);
			\draw[-{Latex[length=3mm, width=2mm]}] (epsD) to (C);
			\draw[-{Latex[length=3mm, width=2mm]}] (B) to (C);
			\draw[-{Latex[length=3mm, width=2mm]}] (epsA) to (A);
			\draw[-{Latex[length=3mm, width=2mm]}] (epsB) to (B);
			\draw[-{Latex[length=3mm, width=2mm]}] (epsC) to (C);
			\draw[{Latex[length=3mm, width=2mm]}-{Latex[length=3mm, width=2mm]}, bend left] (B) to (C);
		\end{tikzpicture}
	\end{center}

	We see that $G(M_{\setminus L})$ is a subgraph of $G(M)^{\setminus L}$.
	\item \label{ans:obs_equiv} The projections of both solution functions $g$ and $\tilde{g}$ onto $V\setminus L = \{A, B, C\}$ are the same, so the result follows directly from the definition of observational equivalence (Definition \ref*{def:scm_obs_interv_equiv}) and Remark \ref*{rem:using_solution_functions} (if the solutions functions are the same, then the Markov kernels coincide).
	\item For both SCMs we have the solution function
	\[\begin{cases}
			g_A(x_B, \varepsilon) = \varepsilon_{A'} \\
			g_C(x_B, \varepsilon) = \varepsilon_{A'} \alpha_{ad} \alpha _{dc}+x_B \alpha _{bc}+\varepsilon_{C'}+\varepsilon_{D'} \alpha_{dc},
	\end{cases}\]
	so similar to \ref{ans:obs_equiv} we conclude that they are observationally equivalent.
\end{enumerate}
\end{gradedsolution}

\end{enumerate}

\end{document}
